\section{Empirical Results}
\label{sec:results}

\subsection{In-Sample Fit}
Table~\ref{tab:insample} details the in-sample HAC-OLS coefficients for the HAR and HAR-S models.
While the HAR-S model improves in-sample $\bar{R}^2$ modestly, the coefficients on the FinBERT
sentiment score and sentiment surprise are statistically insignificant across all regimes for both
SPX and NIFTY (see Section~\ref{sec:regime}). For BTC the primary non-autoregressive in-sample
driver is the squared order-book imbalance (OBI): $p < 0.01$ in the low- and medium-stress
regimes (see Table~\ref{tab:regime_coefs}); FinBERT coefficients for BTC are not identified in
the HAR-S regime splits because the BTC news-sentiment series is unavailable for a sufficient
in-sample window in those partitions.

\begin{table}[htbp]
\centering
\caption{In-Sample HAC-OLS Coefficients: HAR and HAR-S}
\label{tab:insample}
\resizebox{\textwidth}{!}{%
\begin{tabular}{lllrrlrr}
\toprule
Asset & Model & Variable & Coef & p-value & Sig & R\$\textasciicircum \{2\}\$ & N \\
\midrule
BTC & HAR & const & -1.135600 & 0.000000 & *** & 0.403700 & 2744 \\
BTC & HAR & RV\_d & 0.171400 & 0.000000 & *** & 0.403700 & 2744 \\
BTC & HAR & RV\_w & 0.375200 & 0.000000 & *** & 0.403700 & 2744 \\
BTC & HAR & RV\_m & 0.300100 & 0.000000 & *** & 0.403700 & 2744 \\
BTC & HAR-S & const & -1.010300 & 0.124000 & NaN & 0.187300 & 1095 \\
BTC & HAR-S & RV\_d & 0.217500 & 0.000000 & *** & 0.187300 & 1095 \\
BTC & HAR-S & RV\_w & 0.223200 & 0.006000 & *** & 0.187300 & 1095 \\
BTC & HAR-S & RV\_m & 0.538900 & 0.000000 & *** & 0.187300 & 1095 \\
BTC & HAR-S & vpin\_lag1 & -0.794400 & 0.116000 & NaN & 0.187300 & 1095 \\
BTC & HAR-S & obi\_sq\_lag1 & 2.769200 & 0.000000 & *** & 0.187300 & 1095 \\
BTC & HAR-S & illiq\_lag1 & 0.000000 & 0.095000 & * & 0.187300 & 1095 \\
BTC & HAR-S & FinBERT\_score\_lag1 & 0.000000 & 0.252000 & NaN & 0.187300 & 1095 \\
BTC & HAR-S & sent\_surprise\_lag1 & 0.000000 & 0.000000 & *** & 0.187300 & 1095 \\
BTC & HAR-S & credit\_spread\_lag1 & -0.050000 & 0.606000 & NaN & 0.187300 & 1095 \\
BTC & HAR-S & term\_slope\_lag1 & 0.046200 & 0.645000 & NaN & 0.187300 & 1095 \\
BTC & HAR-S & crypto\_fg\_lag1 & 0.020400 & 0.907000 & NaN & 0.187300 & 1095 \\
SPX & HAR & const & -3.360000 & 0.000000 & *** & 0.088500 & 2862 \\
SPX & HAR & RV\_d & -0.004100 & 0.845000 & NaN & 0.088500 & 2862 \\
SPX & HAR & RV\_w & 0.248600 & 0.000000 & *** & 0.088500 & 2862 \\
SPX & HAR & RV\_m & 0.451800 & 0.000000 & *** & 0.088500 & 2862 \\
SPX & HAR-S & const & -9.928200 & 0.000000 & *** & 0.033900 & 755 \\
SPX & HAR-S & RV\_d & -0.038700 & 0.236000 & NaN & 0.033900 & 755 \\
SPX & HAR-S & RV\_w & -0.062100 & 0.549000 & NaN & 0.033900 & 755 \\
SPX & HAR-S & RV\_m & 0.399900 & 0.020000 & ** & 0.033900 & 755 \\
SPX & HAR-S & illiq\_lag1 & -0.000000 & 0.637000 & NaN & 0.033900 & 755 \\
SPX & HAR-S & FinBERT\_score\_lag1 & -0.235200 & 0.802000 & NaN & 0.033900 & 755 \\
SPX & HAR-S & sent\_surprise\_lag1 & -1.084200 & 0.222000 & NaN & 0.033900 & 755 \\
SPX & HAR-S & credit\_spread\_lag1 & 0.628500 & 0.028000 & ** & 0.033900 & 755 \\
SPX & HAR-S & term\_slope\_lag1 & 0.525000 & 0.071000 & * & 0.033900 & 755 \\
NIFTY50 & HAR & const & -4.602000 & 0.000000 & *** & 0.039900 & 2755 \\
NIFTY50 & HAR & RV\_d & -0.016400 & 0.429000 & NaN & 0.039900 & 2755 \\
NIFTY50 & HAR & RV\_w & 0.115200 & 0.034000 & ** & 0.039900 & 2755 \\
NIFTY50 & HAR & RV\_m & 0.482300 & 0.000000 & *** & 0.039900 & 2755 \\
NIFTY50 & HAR-S & const & -10.082400 & 0.000000 & *** & 0.030500 & 719 \\
NIFTY50 & HAR-S & RV\_d & -0.006600 & 0.874000 & NaN & 0.030500 & 719 \\
NIFTY50 & HAR-S & RV\_w & 0.138000 & 0.172000 & NaN & 0.030500 & 719 \\
NIFTY50 & HAR-S & RV\_m & 0.160700 & 0.354000 & NaN & 0.030500 & 719 \\
NIFTY50 & HAR-S & illiq\_lag1 & 0.000000 & 0.978000 & NaN & 0.030500 & 719 \\
NIFTY50 & HAR-S & FinBERT\_score\_lag1 & -0.222700 & 0.803000 & NaN & 0.030500 & 719 \\
NIFTY50 & HAR-S & sent\_surprise\_lag1 & -0.376300 & 0.670000 & NaN & 0.030500 & 719 \\
NIFTY50 & HAR-S & credit\_spread\_lag1 & 0.621900 & 0.022000 & ** & 0.030500 & 719 \\
NIFTY50 & HAR-S & term\_slope\_lag1 & 0.717400 & 0.009000 & *** & 0.030500 & 719 \\
\bottomrule
\end{tabular}
}
\begin{tablenotes}\footnotesize
  \item Significance: $^{*}p<0.10$, $^{**}p<0.05$, $^{***}p<0.01$. HAC standard errors.
\end{tablenotes}
\end{table}


\subsection{Out-of-Sample Forecasting}
\label{sub:oos}

When evaluated out-of-sample, the linear HAR-S model consistently underperforms the baseline HAR
on QLIKE for BTC (DM = $-1.69$, $p = 0.091$) and on both QLIKE and MSE for NIFTY
($\text{QLIKE-DM} = -3.55$, $p < 0.001$). This confirms the well-known tendency for
high-dimensional linear models to overfit noisy sentiment signals out-of-sample. HAR-S does not
improve significantly on HAR for SPX under either loss.  The full-sample OOS metrics for the
baseline models are reported in Table~\ref{tab:oos_full}; the harmonized trailing-window metrics
for all three models are in Table~\ref{tab:oos_metrics}. These results are sourced from
\texttt{table2\_dm\_test.csv} (the canonical DM output of \texttt{run\_paper.py}, which uses
chronologically correct end-aligned array trimming).

The Neural-HAR ensemble shows a different pattern. Table~\ref{tab:dm_tests} (sourced from
\texttt{table2\_dm\_test.csv}) reports QLIKE improvements over the HAR baseline that are
statistically significant at the 5\% level for BTC ($\text{DM} = +1.98$, $p = 0.048$) and SPX
($\text{DM} = +5.05$, $p < 0.001$), and at the 5\% level for NIFTY ($\text{DM} = +2.44$, $p =
0.015$). Under MSE, however, the picture reverses: HAR is preferred over Neural-HAR for all
three assets ($\text{DM}_{\text{MSE}} = -3.28$, $-3.63$, $-3.08$ for BTC, SPX, NIFTY; all
$p < 0.01$). This divergence between QLIKE and MSE results is consistent with the sensitivity
of squared-error loss to large outlier realizations and with \citet{patton2011volatility}'s
argument that QLIKE is the more robust metric when the volatility proxy is itself noisy.

\begin{table}[htbp]
\centering
\caption{OOS Forecast Accuracy: All Models on Harmonised Trailing Window}
\label{tab:oos_metrics}
\resizebox{\textwidth}{!}{%
\begin{tabular}{llrrrr}
\toprule
Asset & Model & N\textbackslash \_OOS & QLIKE & MSE & MAE \\
\midrule
BTC & HAR & 1119 & 0.676900 & 2.034900 & 1.154900 \\
SPX & HAR & 751 & 14.802500 & 7.728000 & 2.118300 \\
NIFTY50 & HAR & 716 & 3.363300 & 6.090500 & 1.810200 \\
\bottomrule
\end{tabular}
}
\end{table}

\begin{table}[htbp]
\centering
\caption{Diebold--Mariano Test (HLN Correction): Challenger vs.\ HAR Benchmark (End-Aligned Trailing Window, $N\approx275$--$287$)}
\label{tab:dm_tests}
\resizebox{\textwidth}{!}{%
\begin{tabular}{llrrllrrllr}
\toprule
Model & Asset & QLIKE\textbackslash \_DM & QLIKE\textbackslash \_p & QLIKE\textbackslash \_sig & QLIKE\textbackslash \_better & MSE\textbackslash \_DM & MSE\textbackslash \_p & MSE\textbackslash \_sig & MSE\textbackslash \_better & OOS\textbackslash \_R2 \\
\midrule
HAR-S & BTC & -1.690300 & 0.091000 & * & A & 1.514900 & 0.129800 & NaN & B & 0.046400 \\
Neural-HAR & BTC & 1.980000 & 0.047700 & ** & B & -3.283300 & 0.001000 & *** & A & -0.152100 \\
HAR-S & SPX & 0.704100 & 0.481300 & NaN & B & 0.082100 & 0.934600 & NaN & B & 0.002200 \\
Neural-HAR & SPX & 5.053000 & 0.000000 & *** & B & -3.630500 & 0.000300 & *** & A & -0.163800 \\
HAR-S & NIFTY & -3.551500 & 0.000400 & *** & A & -2.210700 & 0.027100 & ** & A & -0.049900 \\
Neural-HAR & NIFTY & 2.435800 & 0.014900 & ** & B & -3.077700 & 0.002100 & *** & A & -0.123500 \\
\bottomrule
\end{tabular}
}
\begin{tablenotes}\footnotesize
  \item Positive DM statistic favours the challenger model. `Better' column: A\,=\,HAR preferred, B\,=\,Challenger preferred. All arrays aligned to the chronological end of the shorter series. $^{*}p<0.10$, $^{**}p<0.05$, $^{***}p<0.01$.
\end{tablenotes}
\end{table}


\paragraph{Note on a pipeline artifact in \texttt{table4\_dm\_tests.csv}.}
An earlier supplementary output (\texttt{table4\_dm\_tests.csv}, generated by
\texttt{notebooks/paper\_tables.py}) reported BTC Neural-HAR MSE-DM = $+3.82$ ($p < 0.001$),
the opposite conclusion to the canonical result above. Source-code inspection confirms this was
caused by a front-alignment bug in that notebook: it trimmed all forecast arrays from the start
(\texttt{v[:min\_n]}) rather than the chronological end (\texttt{v[-min\_n:]}), so it was
inadvertently comparing HAR forecasts from the beginning of the HAR series with Neural-HAR
forecasts from the beginning of the Neural-HAR series — two non-overlapping time windows. The
bug has been corrected in the notebook. Throughout this paper we rely exclusively on the output
of \texttt{run\_paper.py}, which has always performed correct end-alignment and whose DM results
are reported in Table~\ref{tab:dm_tests}.

\subsection{Model Confidence Set}
\label{sub:mcs}
Table~\ref{tab:mcs} reports the Model Confidence Set (MCS) results at $\alpha = 0.10$ using
QLIKE losses. For SPX, the Neural-HAR model is the sole survivor; the baseline HAR is eliminated
(MCS $p$-value = 0.0, \texttt{table5\_mcs.csv}). For BTC and NIFTY both HAR and Neural-HAR
survive the MCS, meaning the data do not provide sufficient power to eliminate either model at
the 90\% confidence level over this approximately one-year evaluation window. The full three-model
MCS (including HAR-S) is reported in Table~\ref{tab:mcs_full} for completeness.

\begin{table}[htbp]
\centering
\caption{Model Confidence Set: Neural-HAR vs.\ HAR ($\alpha=0.10$, QLIKE Loss, Trailing Window $N\approx275$--$287$)}
\label{tab:mcs}
\resizebox{\textwidth}{!}{%
\begin{tabular}{llrrrrr}
\toprule
Asset & model & mean\textbackslash \_loss & MCS\textbackslash \_pvalue & in\textbackslash \_MCS & rank & eliminated\textbackslash \_at \\
\midrule
BTC & Neural-HAR & 0.641505 & 1.000000 & True & 1 & NaN \\
BTC & HAR & 0.713312 & 1.000000 & True & 2 & NaN \\
SPX & Neural-HAR & 1.773951 & 1.000000 & True & 1 & NaN \\
SPX & HAR & 5.407347 & 0.000000 & False & 2 & 1.000000 \\
NIFTY50 & Neural-HAR & 1.803965 & 1.000000 & True & 1 & NaN \\
NIFTY50 & HAR & 2.137405 & 1.000000 & True & 2 & NaN \\
\bottomrule
\end{tabular}
}
\end{table}

\begin{table}[htbp]
\centering
\caption{Model Confidence Set: All Three Models ($\alpha=0.10$, QLIKE Loss, End-Aligned Trailing Window)}
\label{tab:mcs_full}
\resizebox{\textwidth}{!}{%
\begin{tabular}{rllrrrrr}
\toprule
Unnamed: 0 & Asset & model & mean\textbackslash \_loss & MCS\textbackslash \_pvalue & in\textbackslash \_MCS & rank & eliminated\textbackslash \_at \\
\midrule
0 & BTC & Neural-HAR & 0.286585 & 1.000000 & True & 1 & NaN \\
1 & BTC & HAR & 0.318403 & 1.000000 & True & 2 & NaN \\
2 & BTC & HAR-S & 0.347834 & 1.000000 & True & 3 & NaN \\
0 & SPX & Neural-HAR & 1.566914 & 1.000000 & True & 1 & NaN \\
1 & SPX & HAR-S & 3.657654 & 0.000000 & False & 2 & 2.000000 \\
2 & SPX & HAR & 3.896581 & 0.000000 & False & 3 & 1.000000 \\
0 & NIFTY & Neural-HAR & 1.869842 & 1.000000 & True & 1 & NaN \\
1 & NIFTY & HAR & 2.330459 & 0.026000 & False & 2 & 2.000000 \\
2 & NIFTY & HAR-S & 3.321181 & 0.009000 & False & 3 & 1.000000 \\
\bottomrule
\end{tabular}
}
\end{table}

